% GENERATED FILE. Edit canonical lessons, then run npm run generate:books.
% canonical-source-hash: fnv1a64:93378af896c401b0
% canonical-lessons: FR-C262-excusez-moi, FR-C262-tout-droit, FR-C262-revision-chemin

\chapter{Finding the Way — excusez-moi, tout droit}
\label{ch:fr-finding-the-way}
\hlchaptermodality{262}

\begin{chapteropening}
\textbf{By the end of this chapter:} \emph{I can stop a stranger politely, ask the way, and follow and give simple directions.}

Complete FR-C262-revision-chemin: ask the way to a place and give the directions back, straight on, left and right.
\end{chapteropening}

\section[excusez-moi, pour aller à… ?]{excusez-moi — the question that finds the way}

\label{lesson:FR-C262-excusez-moi}



\begin{quote}
Say \textquotedblleft{}to go\textquotedblright{}. (\emph{Aller.}) Say \textquotedblleft{}the museum\textquotedblright{}. (\emph{Le musée.})

You are in a town you do not know, and the museum is somewhere.
\end{quote}

\subsection*{You'll want to know: excusez-moi}
\textbf{excusez-moi} — \textquotedblleft{}excuse me\textquotedblright{}, to someone you do not know.

It is the \emph{vous} command of \textbf{excuser}, \textquotedblleft{}to excuse\textquotedblright{}, with \emph{moi} after it: word for word, \textquotedblleft{}excuse me\textquotedblright{}. A stranger in the street is owed \emph{vous}; to a friend it is \textbf{excuse-moi}. The same words apologise: bump into someone and you say \emph{excusez-moi !} The short \textbf{pardon} does both jobs too, and is as polite.

\begin{grammarlens}[title={Grammar Lens: pour aller à…}]
\begin{quote}
\textbf{Excusez-moi, pour aller au musée, s'il vous plaît ?} — Excuse me, how do I get to the museum?
\end{quote}

Word for word, \textquotedblleft{}to go to the museum, please\textquotedblright{}. The rest of the question is left unsaid, and every French speaker hears it anyway. The place takes \textbf{à}, which welds onto \emph{le}:

\begin{quote}
\textbf{au} musée · \textbf{à la} plage · \textbf{à l'}hôtel
\end{quote}

A second way asks for the \textbf{place} rather than the \textbf{way}, with \emph{où est}, \textquotedblleft{}where is\textquotedblright{}. Either can be followed by a short question that saves a walk, \emph{c'est loin}, \textquotedblleft{}is it far\textquotedblright{}, or \emph{c'est près d'ici}, \textquotedblleft{}is it near here\textquotedblright{}:

\begin{quote}
\emph{Où est le musée ? C'est loin ? C'est près d'ici ?}
\end{quote}
\end{grammarlens}

\begin{cousinweb}
\textbf{excuser} is inherited from Latin \emph{excūsāre} (EtymDB 2.1), built on \emph{ex-}, \textquotedblleft{}out of\textquotedblright{}, and \emph{causa}, \textquotedblleft{}a cause, a charge in court\textquotedblright{} (Lewis \& Short, under \emph{excuso}): to release someone from a charge, to free them from blame. So \emph{excusez-moi} asks, at root, to be let off. English \textbf{excuse} is the same word, and \textbf{cause} is its second half.
\end{cousinweb}

\subsection*{Your turn}
Say these aloud:

\begin{itemize}
\raggedright
  \item \emph{Excusez-moi, pour aller au musée, s'il vous plaît ?} — and let the voice rise
  \item the same, to the beach, then to the hotel
  \item ask for the place, and how far — \emph{Excusez-moi, où est le musée ? C'est loin ?}
  \item to a friend — \emph{Excuse-moi !}
\end{itemize}

\begin{tcolorbox}[breakable,colback=teal!4,colframe=teal!35!black,title={Before you move on}]
Who gets \emph{excusez-moi}, and who gets \emph{excuse-moi}? (A \textbf{stranger}; a \textbf{friend}.) Ask the way to the hotel. (\emph{\textbf{Excusez-moi, pour aller à l'hôtel ?}}) What does \emph{à} plus \emph{le} become? (\emph{\textbf{Au.}}) And which Latin word is inside \emph{excuser}? (\emph{\textbf{Causa}}, a charge.)
\end{tcolorbox}
\section[tout droit]{tout droit — straight on, and then the turns}

\label{lesson:FR-C262-tout-droit}



\begin{quote}
Say \textquotedblleft{}on the right\textquotedblright{}. (\emph{À droite.}) Then \textquotedblleft{}on the left\textquotedblright{}. (\emph{À gauche.})

You asked the way, and a stranger is answering. This lesson is what you will hear.
\end{quote}

\subsection*{You'll want to know: tout droit}
\textbf{tout droit} — \textquotedblleft{}straight on\textquotedblright{}.

\begin{quote}
\emph{Allez \textbf{tout droit}.} — Go straight on.
\end{quote}

Listen hard to the end of it. In \textbf{droit} the final \emph{t} is silent: \textquotedblleft{}drwa\textquotedblright{}. In \textbf{droite} it is said: \textquotedblleft{}drwat\textquotedblright{}. \emph{Tout droit} is straight ahead; \emph{à droite} is a turn to the right. The two sound almost alike, and they send you in different directions, so a careful listener waits for that \emph{t}.

\begin{grammarlens}[title={Grammar Lens: directions, as commands}]
A stranger gives directions in the \emph{vous} command: the \emph{vous} form of the present, without the \emph{vous}.

\begin{quote}
\emph{Allez tout droit, puis \textbf{tournez} à gauche.} — Go straight on, then turn left. \emph{\textbf{Continuez} jusqu'au musée.} — Keep going as far as the museum. \emph{\textbf{Traversez} la rue.} — Cross the street. \emph{\textbf{Prenez} la rue à droite.} — Take the street on the right.
\end{quote}

Five verbs carry almost every set of directions: \emph{aller}, \emph{tourner}, \emph{continuer}, \emph{traverser} and \emph{prendre}. The order words come back too: \emph{d'abord}, \emph{puis}, \emph{ensuite}. And a kind stranger may end with a distance: \emph{C'est à deux minutes.} — It's two minutes away.
\end{grammarlens}

\begin{cousinweb}
\textbf{tout} is inherited from Latin \emph{tōtus}, \textquotedblleft{}all, whole\textquotedblright{} (EtymDB 2.1), so \emph{tout droit} is \textquotedblleft{}all straight\textquotedblright{}, the whole way. \textbf{droit} has no parent in EtymDB 2.1; it is usually traced to Latin \emph{dīrēctus}, \textquotedblleft{}set straight\textquotedblright{} (Lewis \& Short, under \emph{directus}), which English kept as \textbf{direct}. \textbf{Droite}, the right hand, is the same word in the feminine. That the right hand was named as the \textquotedblleft{}straight\textquotedblright{} or proper one is the usual account, offered here as that.
\end{cousinweb}

\subsection*{Your turn}
Say these aloud:

\begin{itemize}
\raggedright
  \item \emph{tout droit}, then \emph{à droite} — and hear the \emph{t}
  \item \emph{Allez tout droit, puis tournez à gauche.}
  \item \emph{Continuez jusqu'au musée. Traversez la rue.}
  \item the whole exchange — \emph{Excusez-moi, pour aller à l'hôtel ?} — \emph{Allez tout droit, puis prenez la rue à droite.}
\end{itemize}

\begin{tcolorbox}[breakable,colback=teal!4,colframe=teal!35!black,title={Before you move on}]
How do you hear \emph{tout droit} apart from \emph{à droite}? (The final \emph{\textbf{t}} is said only in \emph{\textbf{droite}}.) Say \textquotedblleft{}turn left\textquotedblright{}. (\emph{\textbf{Tournez à gauche.}}) Which form do the directions take? (The \emph{\textbf{vous}} command.) And what did Latin \emph{tōtus} mean? (\textbf{Whole.})
\end{tcolorbox}
\section[(review)]{Review — finding the way}

\label{lesson:FR-C262-revision-chemin}



\begin{quote}
Stop a stranger politely. (\emph{Excusez-moi.}) Then say \textquotedblleft{}straight on\textquotedblright{}. (\emph{Tout droit.})
\end{quote}

\begin{grammarlens}[title={What you've built}]
Two lessons, and a whole exchange in the street. The opener, \emph{excusez-moi}. The question, \emph{pour aller à… ?}, or \emph{où est… ?} And the answer, in the \emph{vous} command: \emph{allez tout droit}, \emph{tournez à gauche}, \emph{continuez}, \emph{traversez}, \emph{prenez la rue à droite}.

\begin{quote}
— \emph{Excusez-moi, pour aller au musée, s'il vous plaît ?} — \emph{D'abord, allez tout droit. Ensuite, tournez à droite, et traversez la rue. Le musée est là, à côté de l'hôtel.} — \emph{C'est loin ?} — \emph{Non, c'est à deux minutes.}
\end{quote}

The order words from the story chapters are back at work: a set of directions is a little story about the future, told in commands.
\end{grammarlens}

\subsection*{Your turn}
Say these aloud:

\begin{itemize}
\raggedright
  \item the question — \emph{Excusez-moi, pour aller au musée, s'il vous plaît ?}
  \item the answer — \emph{D'abord, allez tout droit. Ensuite, tournez à droite.}
  \item \emph{Traversez la rue. Le musée est à côté de l'hôtel.}
  \item when you did not catch it — \emph{Excusez-moi, vous pouvez m'expliquer encore une fois ?}
  \item directions from your door to the nearest shop, with \emph{d'abord} and \emph{ensuite}
\end{itemize}

\begin{tcolorbox}[breakable,colback=teal!4,colframe=teal!35!black,title={Before you move on}]
What do you listen for to tell \emph{tout droit} from \emph{à droite}? (The final \emph{\textbf{t}} of \emph{droite}.) Say \textquotedblleft{}cross the street\textquotedblright{}. (\emph{\textbf{Traversez la rue.}}) And what do you say to a friend instead of \emph{excusez-moi}? (\emph{\textbf{Excuse-moi.}})
\end{tcolorbox}
